\section{Bounded equations in a semipositive big background}

In this section $\theta$ is again semipositive and big, and we solve the exponential and the normalized equations with $L^p$ densities by bounded potentials. The basic tool is a bounded subsolution. A strict barrier, corrected by the solution of a Hermitian equation, is composed with a bounded convex function; the result is a bounded $\theta$-plurisubharmonic function whose Monge--Amp\`ere measure dominates a prescribed density. The derivative of the convex function decays polynomially near the poles of the barrier, and we compensate for this decay by multiplying the auxiliary density by the reciprocal weight.

Fix a strict barrier $\rho\in\PSH(X,\theta)$ and $0<a\leq1$ such that
\[
 \theta+\ddc\rho\geq a\omega,\qquad \rho\leq-1.
\]
No property of $\rho$ other than these two inequalities will be used.

\Needspace{18\baselineskip}
\begin{lemma}\label{lem:big-lp-subsolution}
Fix $p>1$ and $L>0$. There is $b>0$ such that, for every $f\geq0$ with $\|f\|_{L^p(\mathrm{d}V)}\leq L$, one can find $s\in\PSH(X,\theta)$ satisfying
\[
 -1\leq s\leq0,\qquad \mathrm{MA}_\theta(s)\geq b f\,\mathrm{d}V.
\]
The constant $b$ depends only on the fixed geometry and on $\rho,a,p,L$.
\end{lemma}
\begin{proof}
Choose $1<q<p$. Since a fixed quasi-plurisubharmonic function belongs to every finite $L^r$ space, H\"older's inequality gives a uniform $L^q$ bound for
\[
 g=(-\rho)^{2n}f.
\]
By \cref{thm:kn-continuous} there are $c>0$ and a continuous $\omega$-plurisubharmonic $v$, normalized by $\sup_Xv=-1$, such that
\[
 \omega_v^n=c(1+g)\,\mathrm{d}V.
\]
The construction in that theorem gives a uniform upper bound for $c$, because in its Gauduchon trace estimate the root integral of $1+g$ is bounded below by the root integral of $1$. The lower-density estimate of \cref{thm:lp-potential-estimate} gives
$c\|1+g\|_q\geq c_q>0$, so $c$ is also bounded below. The potential estimate of the same theorem then gives $-C\leq v\leq-1$ with a uniform $C$.

Put $h=\rho+av\leq-1$. Then $\theta_h\geq a\omega_v$. On $(-\infty,-1]$ define
\[
 \chi(t)=-t^{-1}-1,\qquad
 \chi'(t)=t^{-2}\in(0,1],\qquad \chi''(t)=-2t^{-3}>0,
\]
and set $s=\chi(h)$, with value $-1$ at the poles of $h$. To justify both the positivity of $\theta+\ddc s$ and the computation of its Monge--Amp\`ere measure, we first work with the bounded functions $h_k=\max(h,-k)$, $k\geq1$. The chain rule for bounded potentials and $\theta\geq0$ give
\[
 \theta+\ddc\chi(h_k)
 =(1-\chi'(h_k))\theta+\chi'(h_k)\theta_{h_k}
       +\chi''(h_k)\,\mathrm{d}h_k\wedge \dc h_k\geq0.
\]
These bounded potentials decrease to $s$, so $s\in\PSH(X,\theta)$ and $-1\leq s\leq0$. On $\{h>-k\}$, locality and positivity of mixed products give
\[
 \mathrm{MA}_\theta(s)
 \geq a^n|h|^{-2n}\omega_v^n
 \geq a^nc|h|^{-2n}(-\rho)^{2n}f\,\mathrm{d}V.
\]
Letting $k\to\infty$, we obtain this inequality on the set where $h$ is finite; its complement is pluripolar and has zero volume. Since
$|h|\leq(-\rho)+aC\leq(1+aC)(-\rho)$, the assertion follows with any common lower bound for $a^nc(1+aC)^{-2n}$.
\end{proof}

The lemma gives a subsolution for each density in an $L^p$-bounded family. It does not give a lower bound for the masses of all bounded $\theta$-plurisubharmonic functions.

\Needspace{22\baselineskip}
\subsection{The exponential equation}

We approximate $\theta$ from above by the Hermitian forms $\theta+\varepsilon\omega$. The bounded subsolution provides a common lower barrier for the corresponding solutions, and exponential comparison shows that they decrease as $\varepsilon$ decreases. Their limit solves the equation in the original background.

\begin{theorem}\label{thm:big-bounded-exponential}
Let $0\leq f\in L^p(\mathrm{d}V)$, $p>1$, be nonzero, and let $\lambda>0$. There is a unique bounded $v\in\PSH(X,\theta)$ such that
\[
 \mathrm{MA}_\theta(v)=\mathrm{e}^{\lambda v}f\,\mathrm{d}V.
\]
For fixed $\lambda$, these solutions are uniformly bounded when $\|f\|_p$ is bounded above and $\int_X f^{1/n}\,\mathrm{d}V$ is bounded away from zero.
\end{theorem}
\begin{proof}
For $0<\varepsilon\leq1$, set $\theta_\varepsilon=\theta+\varepsilon\omega$. By the Hermitian existence theorem \cref{thm:nguyen-exp}, there is a continuous solution $v_\varepsilon$ of
\[
 \mathrm{MA}_{\theta_\varepsilon}(v_\varepsilon)
       =\mathrm{e}^{\lambda v_\varepsilon}f\,\mathrm{d}V.
\]
Take $s,b$ from \cref{lem:big-lp-subsolution}, and put
$w=s+\lambda^{-1}\log b$. Positivity of bounded products gives
\[
 \mathrm{MA}_{\theta_\varepsilon}(w)
 \geq\mathrm{MA}_\theta(s)\geq bf\,\mathrm{d}V
 \geq \mathrm{e}^{\lambda w}f\,\mathrm{d}V,
\]
and bounded exponential comparison (\cref{cor:bounded-exponential-comparison}) yields $w\leq v_\varepsilon$.
If $0<\varepsilon'<\varepsilon$, then $v_{\varepsilon'}$ is also $\theta_\varepsilon$-plurisubharmonic and
\[
 \mathrm{MA}_{\theta_\varepsilon}(v_{\varepsilon'})
 \geq \mathrm{e}^{\lambda v_{\varepsilon'}}f\,\mathrm{d}V,
\]
so the same comparison gives $v_{\varepsilon'}\leq v_\varepsilon$. Hence $v_\varepsilon$ decreases to a bounded $\theta$-plurisubharmonic function $v$ with $w\leq v\leq v_1$.

To pass to the limit in the equation, we use the fixed Hermitian background $\Omega=\theta+\omega$ and the expansion
\[
 (\theta_\varepsilon+\ddc v_\varepsilon)^n
 =\sum_{k=0}^n\binom nk
       (\Omega+\ddc v_\varepsilon)^k
       \wedge\bigl(-(1-\varepsilon)\omega\bigr)^{n-k}.
\]
Continuity of bounded products along monotone sequences applies to every term, and the smooth coefficients converge. Dominated convergence applies to the right-hand sides. This proves the equation. Uniqueness follows from bounded domination in the semipositive big background (\cref{thm:semipositive-bounded-domination,prop:big-noncollapsing}) by the usual exponential comparison argument.

For the uniform assertion, note first that the lower barrier $w$ is uniform by \cref{lem:big-lp-subsolution}. The upper barrier $v_1$ solves an equation in the fixed Hermitian background $\Omega$. The proof of \cref{lem:exponential-gauduchon-mean} bounds its supremum using only an upper bound for the density norm and a positive lower bound for the root integral: these two bounds control, respectively, the H\"older term and the denominator in \eqref{eq:gauduchon-exponential-mean}, and they survive the smooth approximation used there. Thus $v_1$, and hence $v$, has a uniform upper bound.
\end{proof}

Local continuity of this bounded solution will follow from \cref{thm:local-big-bounded-continuity}.

\subsection{The normalized equation}

For the normalized equation the approximating potentials need not be monotone. Upper envelopes and lower envelopes of their tails provide monotone sequences on either side of an $L^1$ cluster limit. Comparison identifies the two limits, and hence the equation, while the constant is controlled by the Gauduchon trace and the bounded subsolution.

\Needspace{14\baselineskip}
\begin{theorem}\label{thm:big-bounded-normalized}
For every nonzero $0\leq f\in L^p(\mathrm{d}V)$, $p>1$, there are a bounded $\varphi\in\PSH(X,\theta)$ and a uniquely determined $c>0$ such that
\[
 \mathrm{MA}_\theta(\varphi)=cf\,\mathrm{d}V,\qquad \sup_X\varphi=0.
\]
The construction gives uniform bounds for $\|\varphi\|_\infty$, $c$ and $c^{-1}$ when $\|f\|_p$ is bounded above and $\int_Xf^{1/n}\,\mathrm{d}V$ is bounded away from zero.
\end{theorem}

Only the constant $c$ is asserted to be unique; the normalized potential need not be unique, as \cref{thm:degenerate-hermitian-nonuniqueness} shows already for a Hermitian background.

\begin{proof}
\emph{Step 1: approximating equations and their constants.} For integers $j\geq1$, put $\omega_j=\theta+2^{-j}\omega$, and use \cref{thm:kn-continuous}, with the Hermitian form $\omega_j$ in place of $\omega$, to choose continuous normalized solutions
\[
 \mathrm{MA}_{\omega_j}(\varphi_j)=c_jf\,\mathrm{d}V,
 \qquad \sup_X\varphi_j=0.
\]
We first bound $c_j$. Fix a Gauduchon metric $\gamma_0$; when $n=1$ any Hermitian form will do, and the trace estimate reduces to Stokes' theorem. The smooth approximation in the proof of \cref{thm:kn-continuous}, together with the Gauduchon trace, gives a common upper bound: the numerator $\int_X\omega_j\wedge\gamma_0^{n-1}$ is bounded above by $\int_X(\theta+\omega)\wedge\gamma_0^{n-1}$, and the root integral of $f$ is positive and fixed. Here the estimate is applied to the smooth approximating equations before passing to their continuous limits.

For a lower bound, take $s,b$ from \cref{lem:big-lp-subsolution}. If $c_j<b$, then
\[
 \mathrm{MA}_{\omega_j}(\varphi_j)
 \leq(c_j/b)\mathrm{MA}_{\omega_j}(s+A)
\]
for every constant $A$, and strict domination in the Hermitian background $\omega_j$ (\cref{thm:bounded-strict-domination}) would give $\varphi_j\geq s+A$ for every $A$, which is absurd. Thus $c_j\geq b$. In fact the constants decrease with $j$: if $k>j$, positivity gives
$\mathrm{MA}_{\omega_j}(\varphi_k)\geq c_kf\,\mathrm{d}V$, and the same argument excludes $c_j<c_k$. Hence $c_j\downarrow c>0$.

\emph{Step 2: a uniform bound for the potentials.} For the Hermitian form $\Omega=\theta+\omega$, all $\varphi_j$ are normalized $\Omega$-plurisubharmonic functions. Using uniform exponential integrability (\cref{lem:uniform-exponential}) and H\"older's inequality, choose $\gamma>0$ small and $1<q<p$ such that
\[
 g_j=\mathrm{e}^{-\gamma\varphi_j}f
\]
is bounded in $L^q$. Since $\varphi_j\leq0$, we have $g_j\geq f$. \Cref{thm:big-bounded-exponential}, applied with the exponent $q$, therefore gives uniformly bounded $w_j\in\PSH(X,\theta)$ with
\[
 \mathrm{MA}_\theta(w_j)=\mathrm{e}^{\gamma w_j}g_j\,\mathrm{d}V.
\]
Choose $C_0\geq\max\{1,\sup_jc_j\}$ and set $z_j=w_j-\gamma^{-1}\log C_0$. Then
\[
 \mathrm{MA}_{\omega_j}(z_j)
 \geq \mathrm{e}^{\gamma z_j}c_jg_j\,\mathrm{d}V,
 \qquad
 \mathrm{MA}_{\omega_j}(\varphi_j)
       =\mathrm{e}^{\gamma\varphi_j}c_jg_j\,\mathrm{d}V,
\]
and exponential comparison gives $z_j\leq\varphi_j\leq0$. Thus all $\varphi_j$ take values in a common bounded interval.

\emph{Step 3: the upper envelopes.} Extract a subsequence converging in $L^1$ and almost everywhere to the canonical $\theta$-plurisubharmonic representative $\varphi$. We keep the indices $j$; the coefficients of $\omega$ in $\omega_j$ still decrease to zero. Define
\[
 \Phi_j^+=\left(\sup_{\ell\geq j}\varphi_\ell\right)^*,
 \qquad h_j=\inf_{\ell\geq j}\varphi_\ell,
 \qquad p_j=P_\theta(h_j).
\]
All these functions take values in a common bounded interval. The upper envelopes are $\omega_j$-plurisubharmonic and decrease to $\varphi$, since the regularized upper envelope agrees almost everywhere with the $L^1$ limit; this also follows from the local Hartogs lemma applied to the bounded tails.

For $\ell\geq j$, positivity gives $\mathrm{MA}_{\omega_j}(\varphi_\ell)\geq c_\ell f\,\mathrm{d}V\geq cf\,\mathrm{d}V$. By ordered contact (\cref{lem:bounded-ordered-contact}), finite maxima of these functions satisfy the same lower bound. Continuity of bounded products along increasing sequences, followed by the decreasing limit in $j$ and the expansion in the smooth background, gives
\begin{equation}\label{eq:big-normalized-upper-tail}
 \mathrm{MA}_{\omega_j}(\Phi_j^+)\geq cf\,\mathrm{d}V,
 \qquad \mathrm{MA}_\theta(\varphi)\geq cf\,\mathrm{d}V.
\end{equation}

\emph{Step 4: the lower envelopes.} For the reverse inequality we use lower envelopes, which carry a useful exponential weight. For $N\geq j$ put
\[
 p_{j,N}=P_\theta\left(\min_{j\leq\ell\leq N}\varphi_\ell\right).
\]
The obstacle is continuous and bounded, so the measure of $p_{j,N}$ is carried by its contact set. Partition this set into disjoint Borel pieces according to the first index $\ell$ at which $p_{j,N}=\varphi_\ell$. On such a piece, positivity and ordered contact in the background $\omega_\ell$ give
\[
 \mathrm{MA}_\theta(p_{j,N})
 \leq\mathrm{MA}_{\omega_\ell}(p_{j,N})
 \leq c_\ell f\,\mathrm{d}V.
\]
Moreover $p_{j,N}=\varphi_\ell\geq h_j$ there, and $c_\ell\leq c_j$. Summing over the pieces therefore yields the global inequality
\[
 \mathrm{MA}_\theta(p_{j,N})
 \leq \mathrm{e}^{p_{j,N}-h_j}c_jf\,\mathrm{d}V.
\]
As $N\to\infty$, the envelopes decrease to $p_j$. Indeed, their limit lies below each tail obstacle quasi-everywhere and is thus a competitor for $P_\theta(h_j)$, while the reverse inequality is immediate. Continuity of bounded products along monotone sequences and dominated convergence give
\[
 \mathrm{MA}_\theta(p_j)\leq \mathrm{e}^{p_j-h_j}c_jf\,\mathrm{d}V.
\]
The functions $p_j$ increase. Let $\varphi'=(\lim_jp_j)^*$. Since $h_j\to\varphi$ almost everywhere, we have $\varphi'\leq\varphi$ for the canonical representatives. Passing to the increasing limit and using \eqref{eq:big-normalized-upper-tail}, we obtain
\[
 \mathrm{MA}_\theta(\varphi')
 \leq \mathrm{e}^{\varphi'-\varphi}cf\,\mathrm{d}V
 \leq \mathrm{e}^{\varphi'-\varphi}\mathrm{MA}_\theta(\varphi).
\]
Exponential comparison for bounded potentials in the semipositive big background $\theta$, which follows from \cref{thm:semipositive-bounded-domination,prop:big-noncollapsing} as in \cref{cor:bounded-exponential-comparison}, forces $\varphi'\geq\varphi$. Thus $\varphi'=\varphi$, and the first inequality of the last display, together with \eqref{eq:big-normalized-upper-tail}, proves the equation. If necessary, subtracting $\sup_X\varphi$ gives the required normalization.

\emph{Step 5: uniformity and uniqueness of $c$.} All bounds used above depend only on the two stated density bounds and the fixed data, which proves the uniform assertion for the constructed solution. Finally, two different positive constants for the same $f$ would contradict strict domination after an arbitrary translation of one of the bounded solutions, exactly as in the lower bound for $c_j$. This proves the uniqueness of $c$.
\end{proof}

For $\lambda>0$, let $v_\lambda$ be the bounded solution of $\mathrm{MA}_\theta(v_\lambda)=\mathrm{e}^{\lambda v_\lambda}f\,\mathrm{d}V$, and let $(\varphi,c)$ be a normalized solution given by \cref{thm:big-bounded-normalized}. The two exponential barriers used in \cref{prop:vanishing-exponent} give
\[
 \|\lambda v_\lambda-\log c\|_\infty
 \leq\lambda\operatorname{osc}_X\varphi
\]
also in this semipositive background.
